% 第2章 関連研究
% 対応する草稿: research/thesis-drafts/chapter2-related-work-draft.md

\subsection{Confidence Calibration の基礎}\label{subsec:calibration-basics}

% TODO: ECE, Brier Score, AUROC の定義と先行研究

\subsection{LLM の自己評価に関する研究}\label{subsec:llm-self-assessment}

% TODO: Kadavath 2022, Lin 2022, Tian 2023, Xiong 2024

\subsection{動的認識論理（DEL）の基礎}\label{subsec:del-basics}

% TODO: S5/KD45, 知識・信念の公理系
% chapter2-related-work-draft.md §2.3' を LaTeX 化

\subsection{DEL と LLM の先行研究}\label{subsec:del-llm}

% TODO: DEL-ToM (EMNLP 2025), Epistemic Integrity (2024)
% prescriptive vs descriptive の区分

\subsection{日本語 LLM 評価}\label{subsec:japanese-llm}

% TODO: llm-jp-eval, Swallow

\subsection{本研究の位置づけ}\label{subsec:positioning}

% TODO: 2軸の新規性（言語的 + 論理学的）
